\documentclass[11pt, a4paper]{article}

\usepackage[a4paper, top=2.5cm, bottom=2.5cm, left=2cm, right=2cm]{geometry}
\usepackage{amsmath}
\usepackage{amssymb}
\usepackage[colorlinks=true, linkcolor=blue, urlcolor=blue, citecolor=blue]{hyperref}

\title{Theory of Everything - Volume 31 \\ \Large Advanced Game Theory}
\author{Omega Protocol}
\date{\today}

\begin{document}

\maketitle

\section*{Layman's Summary}
Why do people cooperate? Why do economies crash? Volume 31 proves that the mathematical rules governing human behavior and economics are identical to the thermodynamic rules governing quantum networks. Game Theory is just the physics of information flow applied to conscious agents.

\section{Multi-Agent Equilibria (Axiom 34)}
We establish that interacting "rational" nodes—whether they are humans trading stocks, or bacteria competing for food—will mathematically converge to a "Nash Equilibrium." This isn't just behavioral psychology; it is an inescapable consequence of the underlying modular flow seeking out the state of minimum informational resistance.

\section{The Evolution of Cooperation (Theorem)}
In a cutthroat universe, why be nice? We prove that in environments where interactions are repeated (a non-zero-sum game), cooperation isn't just an ethical choice; it's a mathematical guarantee. Strategies like "Tit-for-Tat" naturally emerge because cooperative topological networks dissipate entropy far more efficiently than isolated, purely selfish nodes.

\section{Tragedy of the Commons (Theorem)}
Why do we overfish the oceans and pollute the air? The "Tragedy of the Commons" happens when individuals maximize their own local gain at the expense of the shared environment. We prove that this local optimization, if left completely unchecked by boundary constraints, leads to global topological collapse—a mathematical guarantee of shared ruin.

\end{document}
